Theorem 1. Let $\sigma^{n}=v_{0} v_{1} \ldots v_{n}$ and $\tau^{n}=w_{0} w_{1} \ldots w_{n}$ be simplexes in $\mathbf{R}^{m}$. Then there is a simplicial homeomorphism

$$
\begin{aligned}
& f: \sigma^{n} \leftrightarrow \tau^{n} \\
& f: v_{i} \mapsto w_{i}
\end{aligned}
$$

Proof. For each $v=\sum \alpha_{i} v_{i}\left(\alpha_{i} \geqslant 0, \sum \alpha_{i}=1\right)$, define $f(v)=\sum \alpha_{i} w_{i}$. Then $f$ is bijective, and $f$ and $f^{-1}$ are continuous. (For details, see Problems $0.10-0.17$. Since $f$ and $f^{-1}$ are linear relative to barycentric coordinates, they are linear relative to Cartesian coordinates, and so both are continuous relative to the subspace topology, which we are using, as always.)
